\section{Implicit Function Theorem (1)}

\begin{enumerate}
  \item What problem does the implicit function theorem solve in multivariable calculus?
  \item What does it mean for a function to be defined \emph{implicitly}?
  \item Give an example of an equation that defines a relationship between variables implicitly.
  \item Why can some implicit equations not be solved algebraically for one variable?
  \item What question does the implicit function theorem answer about such equations?
  \item State the implicit function theorem for a function \(F(x,y)\).
  \item What conditions must the function \(F(x,y)\) satisfy for the theorem to apply?
  \item Why is the condition \(F_y(x_0,y_0) \neq 0\) important?
  \item What does the theorem guarantee about the existence of a function \(y=f(x)\)?
  \item In what sense does the resulting function exist (global or local)?
  \item What does ``locally'' mean in the statement of the implicit function theorem?
  \item How does the theorem justify the use of implicit differentiation?
  \item Derive the formula for \(\frac{dy}{dx}\) from the equation \(F(x,y)=0\).
  \item Why does the derivative formula involve the ratio \(-F_x/F_y\)?
  \item What geometric feature often occurs when \(F_y=0\)?
  \item How does the implicit function theorem explain why many curves locally look like graphs of functions?
  \item Give an example of a curve defined implicitly that is not globally a function.
  \item At which points of the circle \(x^2+y^2=1\) does the implicit function theorem fail to produce \(y=f(x)\)?
  \item How can the theorem be used to justify solving equations numerically?
  \item How does the theorem generalize to functions \(F:\mathbb{R}^{n+m} \to \mathbb{R}^m\)?
  \item What role does the Jacobian matrix play in the multivariable version of the theorem?
  \item What rank condition must the Jacobian satisfy?
  \item What does the theorem guarantee about expressing some variables as functions of others?
  \item How many variables can be solved for in terms of the remaining variables?
  \item How does the implicit function theorem relate to the inverse function theorem?
  \item In what sense can the inverse function theorem be viewed as a special case of the implicit function theorem?
  \item How is the implicit function theorem used in constrained optimization?
  \item Why is the theorem important for the method of Lagrange multipliers?
  \item How does the theorem help describe surfaces defined by equations \(F(x,y,z)=0\)?
  \item What geometric object is typically formed by the solution set of such an equation?
  \item What does the theorem imply about the smoothness of these solution sets?
  \item How does the theorem contribute to the concept of a smooth manifold?
  \item Why is the implicit function theorem fundamental in differential geometry?
  \item How does the theorem help justify treating dependent variables as functions in physics and engineering models?
  \item What would go wrong in calculus if the implicit function theorem were not true?
\end{enumerate}
